\documentclass{scrartcl}

\usepackage{wrapfig}
\usepackage{epigraph}
\renewcommand{\epigraphsize}{\scriptsize}
\renewcommand{\epigraphwidth}{60ex}

\usepackage[sexy]{evan}

\hypersetup{
    colorlinks=true,
    linkcolor=blue,
    filecolor=magenta,
    urlcolor=cyan,
    pdftitle={Example},
    pdfpagemode=FullScreen,
}

\urlstyle{same}

\lstset{frame=tb,
  language=Java,
  aboveskip=3mm,
  belowskip=3mm,
  showstringspaces=false,
  columns=flexible,
  basicstyle={\small\ttfamily},
  numbers=none,
  numberstyle=\tiny\color{gray},
  keywordstyle=\color{blue},
  commentstyle=\color{dkgreen},
  stringstyle=\color{mauve},
  breaklines=true,
  breakatwhitespace=true,
  tabsize=3
}

\usepackage{amsmath}
\usepackage{amssymb}

\begin{document}
    \title{INMO 2023/3}
    \author{Reyaansh Agrawal}
    \date{\today}

    \maketitle

    \vspace{0.5cm}

    \begin{problem*}
        Let $\mathbb N$ denote the set of all positive integers. Find all real numbers 
        $c$ for which there exists a function $f:\mathbb N\to \mathbb N$ satisfying:
        \begin{itemize}
            \item for any $x,a\in\mathbb N$, the quantity $\frac{f(x+a)-f(x)}{a}$ is an 
            integer if and only if $a=1$;
            \item for all $x\in \mathbb N$, we have $|f(x)-cx|<2023$.
        \end{itemize}
    \end{problem*}


    \begin{answer*} 
        All \(c=n+\frac 12\) for some nonnegative integer \(n\) work. To verify 
        this, define \(f(n)\) such that it satisfies
        \[f(x)+\frac 14=\begin{cases}cx + \frac 14& x \text{ odd}\\cx + \frac 34& x \text{ even}\end{cases}\]
        This works.
    \end{answer*}
    \begin{soln}
        We will show a claim:

        \begin{claim*} 
            For any \(n\in \mathbb N_0\), we have\[\nu_2(f(x+2^{n})-f(x))=n-1\]
        \end{claim*}
        \begin{proof} 
            Simply induct on \(n\). The base case is clear. Assume this is true for some 
            \(n=k\). Then notice that
            \[n+1 > \nu_2(f(x + 2^{n+1})-f(x))=\nu_2((f(x + 2^{n+1})-f(x+2^n))+(f(x+2^n)-f(x)))>n-1\]
            The claim follows.
        \end{proof}
        Now, notice that for any \(n\in\mathbb{N}, k\in\mathbb{N}_0\) we have\[f(n)\in (cn-2023, cn + 2023)\]and\[f(n + 2^k)\in (c2^k + cn-2023, c2^k + cn + 2023)\]This gives us that\[c2^k\bmod{2^{n-1}}<2\cdot 2023\qquad \text{and}\qquad c2^k\bmod 2^k \ge 2\cdot 2023\]Equivalently,\[\{2c\}=2c\bmod 1<\frac{2023}{2^{k-2}}\qquad\text{and}\qquad \{c\}=c\bmod 1 \ge \frac{2023}{2^{k-1}}\]Evidently, as \(n\rightarrow \infty\), we get \(\{2c\}=0\). The second inequality rules out \(\{c\}=0\), giving us the claimed solutions.
    \end{soln}
\end{document}